% (User input window)

\paragraph{Schur--Dirichlet 扭点因子化：通道 torsion 取值表的 Dirichlet--Fourier 分裂}\label{subsubsec:pom-schur-dirichlet-torsion-factorization}
定理 \ref{thm:pom-schur-cycle-index-tomography} 给出 \(S_q\)-等变 Schur 分解
\[
A_q\ \cong\ \bigoplus_{\lambda\vdash q}\bigl(B_{q,\lambda}\otimes I_{f^\lambda}\bigr),
\qquad f^\lambda:=\dim S^\lambda,
\]
并将每个 Schur 通道的迹序列编码为特征标加权的有限组合。
本段补充一个与有限部“primitive 扭点表 \(\Rightarrow\) 角色分辨幂和”平行的接口：对每个 Schur 通道，在 primitive \(q\)-次单位根处取值并做 Dirichlet 角色分解，可得到一组角色扭点泛函，它们在解析上等价于该通道的归一化谱幂和切片。

\paragraph{设定：归一化扭点多项式与角色扭点泛函}
固定整数 \(q\ge2\)，记 \(U(q):=(\ZZ/q\ZZ)^\times\) 与 \(\omega_q:=e^{2\pi i/q}\)，并记平凡通道的 Perron 根
\[
\rho_q:=\rho(B_{q,(q)}).
\]
对每个通道 \(\lambda\vdash q\) 定义归一化扭点多项式
\[
F_{q,\lambda}(t)
:=
\frac{\det\!\bigl(I-tB_{q,\lambda}/\rho_q\bigr)}
{(1-t)^{\mathbf 1_{\{\lambda=(q)\}}}},
\]
并定义归一化（去平凡因子）幂和
\[
u_{n,\lambda}
:=
\frac{\Tr(B_{q,\lambda}^{\,n})}{\rho_q^{\,n}}
-\mathbf 1_{\{\lambda=(q)\}}
\qquad(n\ge1).
\]
于是 \(F_{q,\lambda}(t)=\prod_{\mu\in\mathrm{spec}'(B_{q,\lambda})}(1-(\mu/\rho_q)t)\)，其中 \(\mathrm{spec}'\) 表示当 \(\lambda=(q)\) 时删去 Perron 根 \(\rho_q\)（按重数），否则取全谱。
特别地，因 \(|\mu/\rho_q|<1\) 对所有 \(\mu\in\mathrm{spec}'(B_{q,\lambda})\) 成立，故 \(\log F_{q,\lambda}\) 在 \(|t|\le1\) 上可由 \(F_{q,\lambda}(0)=1\) 的解析分支一致确定。

\begin{definition}[Schur--Dirichlet 角色扭点泛函]\label{def:pom-schur-dirichlet-torsion-functional}
对任意 Dirichlet 角色 \(\chi\)（模 \(q\)，允许非本原，视作 \(U(q)\) 的角色）定义
\[
\mathcal{L}_{q,\lambda}(\chi)
:=
-\sum_{a\in U(q)}\chi(a)\,\log F_{q,\lambda}(\omega_q^{\,a})
\ \in\ \CC/2\pi i\ZZ.
\]
\end{definition}

\begin{theorem}[Schur--Dirichlet 因子化与 Gauss--Dirichlet 展开]\label{thm:pom-schur-dirichlet-factorization}
\leanverified{paper\_pom\_schur\_dirichlet\_factorization}
令
\[
G(\chi;n):=\sum_{a\in U(q)}\chi(a)\,\omega_q^{\,an},
\qquad
\tau(\chi):=G(\chi;1)=\sum_{a\in U(q)}\chi(a)\omega_q^{\,a}.
\]
则对任意 \(\lambda\vdash q\) 与任意 \(\chi\) 有绝对收敛恒等式
\[
\boxed{\;
\mathcal{L}_{q,\lambda}(\chi)
=
\sum_{n\ge1}\frac{u_{n,\lambda}}{n}\,G(\chi;n)\; }.
\]
并且当 \((n,q)=1\) 时有 \(G(\chi;n)=\overline{\chi}(n)\tau(\chi)\)。
若 \(\chi\) 为本原角色（导子为 \(q\)），则 \(\tau(\chi)\neq0\)，且以 Dirichlet 约定 \(\chi(n)=0\)（当 \((n,q)>1\)）时，
\[
\boxed{\;
\mathcal{L}_{q,\lambda}(\chi)
=
\tau(\chi)\sum_{\substack{n\ge1\\(n,q)=1}}\frac{u_{n,\lambda}}{n}\,\overline{\chi}(n)\; }.
\]
此外，\(\{\mathcal{L}_{q,\lambda}(\chi)\}_{\chi\in\widehat{U(q)}}\) 与 primitive torsion 表 \(\{\log F_{q,\lambda}(\omega_q^{\,a})\}_{a\in U(q)}\) 在 Dirichlet 角色变换下互为精确反演：
\[
\boxed{\;
-\log F_{q,\lambda}(\omega_q^{\,a})
=
\frac{1}{\varphi(q)}\sum_{\chi\in\widehat{U(q)}}\overline{\chi}(a)\,\mathcal{L}_{q,\lambda}(\chi)
\qquad(a\in U(q))\; }.
\]
\end{theorem}
\begin{proof}
对任意 \(|x|<1\) 有 \(-\log(1-x)=\sum_{n\ge1}x^n/n\)。
由 \(F_{q,\lambda}(t)=\prod_{\mu\in\mathrm{spec}'(B_{q,\lambda})}(1-(\mu/\rho_q)t)\) 与 \(|\mu/\rho_q|<1\) 得
\[
-\log F_{q,\lambda}(\omega_q^{\,a})
=
\sum_{n\ge1}\frac{u_{n,\lambda}}{n}\,\omega_q^{\,an},
\]
且该级数绝对收敛允许逐项求和。
将其乘以 \(\chi(a)\) 并对 \(a\in U(q)\) 求和，交换求和次序即得第一条恒等式。
当 \((n,q)=1\) 时，乘法 \(a\mapsto an\) 为 \(U(q)\) 上置换，从而 \(G(\chi;n)=\overline{\chi}(n)\tau(\chi)\)；
本原情形下对 \((n,q)>1\) 有 \(G(\chi;n)=0\) 为标准 Gauss 和性质。
最后一式为 \(U(q)\) 上的角色正交反演：对函数
\(
f(a):=-\log F_{q,\lambda}(\omega_q^{\,a})
\)
有 \(\mathcal{L}_{q,\lambda}(\chi)=\sum_{a\in U(q)}\chi(a)f(a)\)，故
\(
f(a)=\varphi(q)^{-1}\sum_{\chi}\overline{\chi}(a)\mathcal{L}_{q,\lambda}(\chi)
\)。
证毕。
\end{proof}

\begin{theorem}[“角色--Schur--torsion”三重接口：谱包络、常数异常与 prime-shadow 证书]\label{thm:pom-schur-dirichlet-torsion-triple-interface}
固定整数 \(q\ge2\) 并沿用上文记号。
假设对每个通道 \(\lambda\vdash q\)，归一化扭点多项式 \(F_{q,\lambda}(t)\) 的次数满足
\[
\deg F_{q,\lambda}\ \le\ \varphi(q),
\]
从而 primitive torsion 表 \(\{F_{q,\lambda}(\omega_q^{\,a})\}_{a\in U(q)}\) 在代数上可唯一插值恢复 \(F_{q,\lambda}\)。
则对固定 \(q\)，有限数据
\[
\Bigl\{\mathcal{L}_{q,\lambda}(\chi)\Bigr\}_{\lambda\vdash q,\ \chi\in\widehat{U(q)}}
\quad\text{与}\quad
\mathcal{M}_{q,(q)}:=\det\nolimits'\!\bigl(I-B_{q,(q)}/\rho_q\bigr)^{-1}
\]
唯一决定以下三类对象：
\begin{enumerate}
  \item \textbf{（指数级）}各通道的归一化谱包络 \(\max_{\mu\in\mathrm{spec}'(B_{q,\lambda})}|\mu|/\rho_q\)，从而在给定 \(\rho_q\) 的标度下确定 \(\rho(B_{q,\lambda})\)，并确定 near-RH 缺陷
  \(
  \Delta_q=\max_{\lambda\neq(q)}\bigl(2\log\rho(B_{q,\lambda})-\log\rho_q\bigr)
  \)
  （见 \S\ref{subsubsec:pom-wreath-artin-factorization} 的定义）；
  \item \textbf{（常数级）}各通道的常数层异常
  \[
  \mathcal{M}_{q,\lambda}:=
  \begin{cases}
  \mathcal{M}_{q,(q)}, & \lambda=(q),\\[2pt]
  \det\!\bigl(I-B_{q,\lambda}/\rho_q\bigr)^{-1}=F_{q,\lambda}(1)^{-1}, & \lambda\neq(q),
  \end{cases}
  \]
  以及由这些常数的有限线性组合所生成的总体 moment-anomaly；
  \item \textbf{（算术级）}Schur prime-shadow 的平方根门槛证书：将各通道的 \(\rho(B_{q,\lambda})\) 与常数输入定理 \ref{thm:pom-schur-chebotarev-prime-deviation}，可得到对通道 prime-偏差的全部可核验上界，且 near-RH 过临界事件 \(\Delta_q>0\) 与某一非平凡通道偏差穿越平方根门槛等价（同定理 \ref{thm:pom-schur-chebotarev-prime-deviation}）。
\end{enumerate}
从而，对固定 \(q\) 而言，“角色扭点表”同时是：Schur 通道谱的断层扫描、常数级异常的可加坐标、以及 prime-shadow 偏差的解析证书；三者由同一份有限数据在三条正交投影下得到。
\end{theorem}
\begin{proof}[证明要点]
由定理 \ref{thm:pom-schur-dirichlet-factorization} 的反演式，\(\{\mathcal{L}_{q,\lambda}(\chi)\}_\chi\) 与 primitive torsion 表互为等价数据；
在次数界 \(\deg F_{q,\lambda}\le\varphi(q)\) 下，由插值唯一性可恢复 \(F_{q,\lambda}\) 的系数与根集，从而恢复归一化谱包络并得到各常数 \(\mathcal{M}_{q,\lambda}=F_{q,\lambda}(1)^{-1}\)。
最后一条 prime-shadow 证书由定理 \ref{thm:pom-schur-chebotarev-prime-deviation} 直接给出。证毕。
\end{proof}

\begin{remark}[扭点对数分支的规范化与代数可核验性]\label{rem:pom-schur-dirichlet-log-branch-normalization}
由于对所有 \(\mu\in\mathrm{spec}'(B_{q,\lambda})\) 有 \(|\mu/\rho_q|<1\)，归一化扭点多项式
\(
F_{q,\lambda}(t)=\prod_{\mu\in\mathrm{spec}'(B_{q,\lambda})}(1-(\mu/\rho_q)t)
\)
在闭单位盘 \(\{|t|\le1\}\) 内无零点，
从而存在唯一全纯对数 \(\log F_{q,\lambda}\) 满足 \(\log F_{q,\lambda}(0)=0\)；
因此 \(\log F_{q,\lambda}(\omega_q^{\,a})\) 与 \(\mathcal{L}_{q,\lambda}(\chi)\) 在本文制度下均为规范确定的复数，
写作 \(\CC/2\pi i\ZZ\) 仅用于强调分支不变性（与有限部的 remark \ref{rem:finite-part-torsion-log-branch-normalization} 同型）。
\par
若进一步 \(F_{q,\lambda}(t)\) 的系数落在某个数域 \(K\)（例如由 \(B_{q,\lambda}\) 的特征多项式所生成），则 primitive torsion 表
\(\{F_{q,\lambda}(\omega_q^{\,a})\}_{a\in U(q)}\) 属于 \(K(\omega_q)\)，
从而可在 Galois 作用下与其 Dirichlet 角色分解一致追踪（对数层面通过上述规范化消除全局 \(2\pi i\ZZ\) 不定性）。
\end{remark}
